\begin{abcliste}
\abc
Falsch ist \result{E_\mathrm{pot} = \EwP{0}{3} \text{ in (1)}}. Die Schülerin hat die Höhe $h$ über dem oberen Ende der Feder genommen, also über (2), obwohl sie (3) als Nullniveau gewählt hat. Von (3) aus liegt (1) um $h + s$ höher:
\begin{align}
 E_\mathrm{pot,1} &= \EaF = \mB\cdot\g\cdot\left(\hB+\sB\right) = \Ea \approx \result{\EaP{0}{3}}
\end{align}
Die anderen Einträge stimmen: In (2) ist die Kugel $s$ über (3), $m\,g\,s = \EbP{0}{3}$, und ihre kinetische Energie ist $m\,g\,h$, weil sie von (1) bis (2) um $h$ gefallen ist. In (3) ist sie in Ruhe, und die ganze Energie steckt in der Feder: $\frac{1}{2}\,k\,s^2 = \EcP{0}{3}$.

\abc
Ohne Reibung muss die Summe in jeder Zeile gleich sein. Nach der Tabelle ist sie in (1) nur \EwP{0}{3}, in (2) aber $\EbP{0}{3} + \EwP{0}{3} = \EaP{0}{3}$ und in (3) \EcP{0}{3}: Die Energie würde zunehmen. Mit dem korrigierten Eintrag ist die Summe überall \EaP{0}{3}. Beim Zusammendrücken sinkt die Kugel nämlich noch um $s$ weiter; auch diese Lageenergie $m\,g\,s$ geht in die Feder.

\abc
Von (1) bis (2) fällt sie frei um $h$:
\begin{align}
 v_2 &= \vbF = \sqrt{2\cdot\g\cdot\hB} = \vb \approx \result{\vbP{0}{2}}
\end{align}
\end{abcliste}
